\section{A free-scale carrier of degree 44}
\label{sec:free-scale15}
This endpoint uses fifteen unrestricted integer variables, has exact degree
$44$, and a division-free evaluator with $111$ binary arithmetic gates
($63$ multiplications, $26$ additions, $22$ subtractions).
It improves the prior degree $48$ at a cost of three gates; it does not
Pareto-dominate the retained $108$-gate circuit. The same elementary
projection factorization reduces that old evaluator to $107$ gates,
so the apples-to-apples evaluation cost is four gates.
The new device is to release the scale variable, replace premature
coprimality by a magnitude guard, and recover the exact scale only after
an amplified whole-product projection.

\subsection{Complete polynomial}
The variables are
\[
 j,C,z,h,r,v,w,T,c,n,f,y,p_1,p_2,p_3\in\mathbb Z.
\]
The fixed numerals $\beta,K_0,M_C,M_F,A_c,B_c$ are exactly the six
pinned source34 coefficient payloads; $M_F$ is already shifted.
All the following are polynomial abbreviations, not witnesses:
\begin{gather*}
 J=6j+1,\quad q_m=\beta J,\quad q=q_m+1,\quad W=12w+4,\\
 X=qW,\quad Y=v,\quad F=(K_0+W)C-zq_m,\\
 R=(q(q-F)-C)(q^2-1)+(M_C+qM_F)J,\\
 E=XY,\quad a=E+Y,\quad H=4a+3,\quad \Delta=a^2+H,\\
 k=R+1+hE,\quad \eta=c-Yk,\quad L=E(Yk),\\
 I=(2r+1)c^2,\quad Q=\Delta(f^2-1),\quad V=f^2(cT-R)-c,\\
 M=2q^3,\quad U=MH,\quad K_{\rm lo}=2X^2-3Xc-2,\\
 Z=Un+MK_{\rm lo}+H(Y-q^3),\\
 D_{\rm err}=(A_cF-q)^2+(B_cC-q)^2,\\
 B_0=k^2(q^2-D_{\rm err})-4q^2\eta^2,\\
 B_{\rm new}=Y^2B_0-k^2q^8,\qquad
 P_{\rm os}=1+p_1^2+p_2^2+p_3(p_3+1).
\end{gather*}
Define the residuals
\begin{align*}
 A&=(Z+Uk)^2+4U(LZ-U),\\
 B&=f^2-\Delta I^2-1,\\
 C_{\rm aux}&=V^2-(Q-1)y(2V+y)-1,\\
 G&=B_{\rm new}-P_{\rm os}.
\end{align*}
The evaluator uses the equivalent schedules
\[
 K_{\rm lo}=X(2X-3c)-2,\qquad
 B_{\rm new}=(Yk)^2(q^2-D_{\rm err})
 -q^2\bigl(4(Y\eta)^2+(kq^3)^2\bigr),
\]
reusing $Yk$ and $q^3$. The sole equation is
\[
 \boxed{P=A^2+B^2+C_{\rm aux}^2+G^2=0.}
\]

\subsection{Monicity before coprimality}
The actual $\beta=2^{1064921}-1$ is odd and $1\pmod3$.
Consequently $q$ is nonzero, even, and $2\pmod3$.
The polynomial $H$ is odd and nonzero, so $U\ne0$ on every assignment.
At a zero, $t=Z/U\in\mathbb Q$ satisfies
\[
 t^2+2(2L+k)t+k^2-4=0.
\]
Monicity makes $t$ integral. Modulo $2$, $t+k$ is even, so
$s=(t+k)/2$ is integral and
\[
 s^2+(2s-k)L-1=0,\qquad (s+L)^2=1+L(L+k).
\]
The same divisibility $U\mid Z$ gives
\[
 M\mid H(Y-q^3),\qquad H\mid MK_{\rm lo}.
\]
There is deliberately no early assertion that $H$ and $q$ are coprime.
Since $H$ is odd and $q$ even, the first relation nevertheless gives
$Y$ even, hence $\Delta$ odd. If
$K=X(X+2ac)-1$, then $2K-K_{\rm lo}=HXc$, so oddness of $H$ gives
\[
 H\mid q^3 K.
\]

\subsection{The guard replaces the early scale hypothesis}
The polynomial $P_{\rm os}$ is strictly positive and represents every
positive integer by the three-square theorem. Hence
\[
 Y^2B_0>k^2q^8>0,
\]
where $k\ne0$ follows because $k=0$ would make $B_{\rm new}\le0$.
Since $B_0\le k^2q^2$, this proves
\[
 B_0>0,\qquad |Y|>|q|^3.
\]
The unchanged centered-budget argument gives common strict signs for
$q,F,C,\alpha=q-F-C$, and
\[
 |\eta/k|<\tfrac12,\qquad
 |k|(|Y|-\tfrac12)<|c|<|k|(|Y|+\tfrac12),\qquad
 \operatorname{sgn}(c)=\operatorname{sgn}(kY).
\]
Put $q_*=|q|$, $x=|X|$, $m=|Y|$. The authentic packing estimates give
\[
 q_*\ge16,\quad q_*(q_*^2-1)<R<q_*^4-2,\quad
 x\ge4q_*\ge64,\quad |E|>4q_*^4>3R+1.
\]
These replace the old literal cubic-scale expression wherever it was
used before projection. Its remaining parity consequence, $\Delta$ odd,
was already proved without coprimality.

\subsection{Signed Pell bootstrap: inherited statements and hypotheses}
The unchanged strong and auxiliary residuals recover the missing main
Pell ordinate by the predecessor's odd-index folding argument.
For $B_p=|a+2|$, $c_*=|c|$, and $I_*=|2r+1|c_*^2$, they give
\[
 |f|^2-(B_p^2-1)I_*^2=1,\qquad |f|>2I_*.
\]
With $S=\Delta I_*$ the auxiliary row is
\[
 S^2V^2-(S^2-1)(V+y)^2=1.
\]
The auxiliary Pell index is odd. Its polynomial identity, reduced
modulo $|f|$, gives an ordinate congruent to $\pm c_*$.
Folding the index into the interval from $0$ to the strong index
makes the ordinate at most $I_*$. Since $I_*+c_*<|f|$, this
congruence gives an actual main ordinate $c_*=\psi_{B_p}(p)$.

The first Pell classification gives
$|k|=2\psi_{|2XY^2+1|}(n_0)$.
The inequality $|c|>|k|$, the parameter comparison, the congruence
$k=R+1+hE$, and $|E|>3R+1$ give the same signed first-index lower
bound as in the source. The unchanged strong-rank lemma and signed
auxiliary step-down then give
\[
 p=R,\qquad 2n_0=R+1,\qquad R=2s_0+1,\quad s_0\ge24.
\]
The sign-uniform ratio estimate uses only these classifications,
$x\ge64$, $m\ge6$, and $|c/k|<m+1$. It proves
\[
 m>\frac{(x-3)^{s_0}}2-1,\qquad |H|>x(x-3)^{s_0}.
\]
These are precisely the inherited unbounded Pell steps, not conclusions
of finite testing. Their complete statements and the verified hypotheses
are detailed in \texttt{PROOF.md}, Sections 3--4, with pinned source texts.
No main projection or power-of-two property has yet been used.

\subsection{Amplified projection and late scale recovery}
Let $\sigma=\operatorname{sgn}(c)$ and $x_0=\sigma X$.
The signed polynomial Pell projections give
\[
 2^R K\equiv(x_0-2^R)(2^Rx_0+1)\pmod H.
\]
Thus $H$ divides $q^3$ times this whole product. We bound that
larger product instead of canceling a possible nonunit modulo $H$.
For $x\ge64$ and $s_0\ge24$,
\[
 (x/32)^{s_0}\ge(x/32)^{24}\ge x^4,
\]
since $x^{20}\ge64^{20}=32^{24}$. Therefore
\[
 |H|>x(x-3)^{s_0}\ge x(x/2)^{s_0}\ge16^{s_0}x^5.
\]
Since $q_*\le x/4$ and $R=2s_0+1$,
\[
 q_*^3 2^{2R+2}x^2
 \le\frac{x^3}{64}16^{s_0+1}x^2
 =\tfrac14 16^{s_0}x^5<|H|.
\]
This strictly bounds the absolute value of the divisible whole product.
The product is zero. Its factor $2^Rx_0+1$ is odd and cannot vanish,
so $x_0=2^R$. Because $X\equiv2\pmod3$ and $R$ is odd,
$\sigma=1$ and $X=2^R$.

The zero whole product and the same Pell congruence now give
$H\mid2^RK$. Since $H$ is odd, $H\mid K$ and hence $H\mid K_{\rm lo}$.
Set $g=K_{\rm lo}/H\in\mathbb Z$. Substituting $K_{\rm lo}=Hg$ into
$MH\mid Z$ and canceling the nonzero integer $H$ gives
\[
 M\mid Y-q^3,\qquad v_{\rm scale}=(Y-q^3)/(2q^3)\in\mathbb Z.
\]
This is cancellation in the integers, not cancellation of a nonunit
modulo $H$. In particular, no CRT coprimality claim is needed.
Thus the exact old cubic scale is recovered after, rather than before,
the signed bootstrap. The surviving index sign gives $k>0$; with
$c>0$, this also gives $Y>0$. The inherited centered half-binomial
extraction now applies without change.

\subsection{Explicit maps and canonical completeness}
From any new zero construct a degree-$48$ predecessor zero as follows.
Retain $j,C,z,h,r,w,T,c,f,y$, and set
\[
 v_{\rm old}=Y/q,\qquad n_{\rm old}=t-g-v_{\rm scale}.
\]
The old carrier's quotient is exactly
\[
 \frac{Z_{\rm old}}{U_{\rm old}}
 =n_{\rm old}+g+\frac{v_{\rm old}-q^2}{2q^2}=t.
\]
Its strong and auxiliary rows are unchanged. Represent the positive
integer $B_0$ by its three Pos coordinates. This restores a literal
old zero, so the old native/compiler/refutation theorem applies.

Conversely start with a finite syntactic refutation of ZFC and use the
same authenticated canonical compiler family. It supplies an old zero
with $q>0$, $X=2^R$, $s_0\ge24$, the exact cubic scale, and
\[
 D_{\rm err}\le q^2/2,\quad 0<\eta/k<1/4,\quad
 B_0>q^2k^2/4,\quad Y\ge X^{s_0}/2>2q^3.
\]
Consequently
\[
 B_{\rm new}>k^2q^2(Y^2/4-q^6)>0.
\]
Retain the arithmetic coordinates, set $v=Y$ and
$n=t-g-v_{\rm scale}$, and represent $B_{\rm new}$ by the three Pos
coordinates. Then $Z/U=t$, and all four new rows vanish.
These constructions prove equivalence of zero existence; a bijection
of arbitrary zero sets is not asserted.

\subsection{Exact degree, reproducibility, and boundary}
The four residual degrees are $22,12,16,18$. In degree $22$, the
first residual contributes only $4ULZ$. With subscript $0$ for
leading homogeneous forms,
\begin{gather*}
 q_0=6\beta j,\quad X_0=12q_0w,\quad Y_0=v,\\
 R_0=-q_0^3(12wC-q_0z),\quad
 U_0=8q_0^3X_0Y_0,\quad L_0=X_0Y_0^2R_0,\\
 Z_0=U_0n+4q_0^3X_0^2,\qquad P_{44}=16U_0^2L_0^2Z_0^2.
\end{gather*}
In particular,
\[
 [j^{26}w^6v^8z^2n^2]P=2^{54}3^{32}\beta^{26}>0.
\]
The exponents sum to $44$, and the other squared rows have smaller
degree. The exact degree therefore holds after fixing every coefficient.

The accompanying independent cut-register audit authenticates coefficient
bytes, checks the literal graph, exact identities and leading coefficient,
and rejects altered arithmetic gates. These finite checks validate the
transcription and algebra, not the unbounded inherited compiler theorem.
The result is an ordinary arithmetic equivalence, hence an equivalence
provable in ordinary ZFC between no integer zeros and syntactic
$\operatorname{Con}(\mathrm{ZFC})$. It is not a proof of consistency,
a global optimum, or a new formal-kernel certificate.
